\documentclass[a4,11pt]{aleph-notas}

% -- Paquetes adicionales
\usepackage{enumitem}
\usepackage{aleph-comandos}
\usepackage{systeme}

% -- Datos
\institucion{Facultad de Ciencias Exactas, Naturales y Ambientales}
\carrera{Catálogo STEM}
\asignatura{Álgebra lineal}
\tema{Ejercicios 09: Espacios con producto interno}
\autor[A. Merino]{Andrés Merino}
\fecha{Periodo 2026-1} 

\logouno[0.16\textwidth]{Logos/logoPUCE_04_ac}
\definecolor{colortext}{HTML}{1957d1}
\definecolor{colordef}{HTML}{1957d1}
\fuente{montserrat}

\begin{document}

\encabezado


%%%%%%%%%%%%%%%%%%%%%%%%%%%%%%%%%%%%%%
\section{Espacios con producto interno}
%%%%%%%%%%%%%%%%%%%%%%%%%%%%%%%%%%%%%%

%%%%%%%%%%%%%%%%%%%%%%%%%%%%%%%%%%%%%%%%
\begin{ejer}
    En cada espacio vectorial $V$ utilice el proceso de Gram-Schmidt para transformar la base $B$ de $V$ en (a) una base ortogonal; (b) una base ortonormal.
    \begin{enumerate}
        \item En $V=\R^{2}$ con $B=\{(1, 2), (-3, 4)\}$.
        \item En $V=\R^{3}$ con $B=\{(1, 1, 1), (0, 1, 1), (1, 2, 3)\}$.
    \end{enumerate}
\end{ejer}

\begin{proof}[Solución]\hspace{0pt}
    \begin{enumerate}
    \item Sean $u_{1} = (1, 2)$ y $u_{2} = (-3, 4)$, vamos a determinar la base ortogonal $B_{1}=\{v_{1}, v_{2}\}$, para hacerlo utilizamos el proceso de Gram-Schmidt.
    \begin{itemize}
        \item Para determinar $v_{1}$: 
        \[
            v_{1} = u_{1} = (1, 2).
        \]
        \item Para determinar $v_{2}$:
        \begin{align*}
            v_{2} &= u_{2} - \dfrac{u_{2} \cdot v_{1}}{v_{1} \cdot v_{1}}v_{1}\\
                  &= (-3, 4) - \dfrac{(-3, 4) \cdot (1, 2)}{(1, 2) \cdot (1, 2)}(1, 2)\\
                  &= (-3, 4) - (1, 2)\\
                  &= (-4, 2).
        \end{align*}
    \end{itemize}
    Vamos a determinar ahora la correspondiente base ortonormal $B_{2}=\{w_{1}, w_{2}\}$. Para esto, notemos que
    \begin{align*}
        w_{1} &= \dfrac{1}{\|v_{1}\|}v_{1} = \dfrac{1}{\sqrt{5}}(1, 2).\\
        w_{2} &= \dfrac{1}{\|v_{2}\|}v_{2} = \dfrac{1}{2\sqrt{5}}(-4, 2).
    \end{align*}
    Por lo tanto, $B_{2}=\left\{\left(\dfrac{\sqrt{5}}{5}, \dfrac{2\sqrt{5}}{5}\right), \left(-\dfrac{2\sqrt{5}}{5}, \dfrac{\sqrt{5}}{5}\right)\right\}$.
    \item Verifique que la base ortogonal es
    \[
        B_{1} = \{(1, 1, 1), (-2, 1, 1), (0, -1, 1)\}
    \]
    y que la base ortonormal es
    \[
        B_{2} = \left\{\left(\dfrac{\sqrt{3}}{3}, \dfrac{\sqrt{3}}{3}, \dfrac{\sqrt{3}}{3}\right), \left(-\dfrac{\sqrt{6}}{3}, \dfrac{\sqrt{6}}{6}, \dfrac{\sqrt{6}}{6}\right), \left(0, -\dfrac{\sqrt{2}}{2}, \dfrac{\sqrt{2}}{2}\right)\right\}. \qedhere
    \]
\end{enumerate}
\end{proof}


%%%%%%%%%%%%%%%%%%%%%%%%%%%%%%%%%%%%%%%%
\begin{ejer}
    En cada caso, calcule el producto interno indicado.
    \begin{enumerate}
        \item En $\R^4$ con el producto interno usual:
        \[
            x=(1,-2,0,3),\quad y=(4,1,-5,2).
        \]
        \item En $\Mat[\R]{2}{2}$ con
        \[
            \langle A,B\rangle = \tr(AB^\intercal),
        \]
        para
        \[
            A=\begin{pmatrix}1&2\\-1&0\end{pmatrix},\quad
            B=\begin{pmatrix}3&-1\\4&2\end{pmatrix}.
        \]
        \item En $\R_2[x]$ con
        \[
            \langle p,q\rangle = a_0b_0+a_1b_1+a_2b_2,
        \]
        donde
        \[
            p(x)=2-x+3x^2,\quad q(x)=1+4x-x^2.
        \]
        \item En $\mathcal{C}([0,1])$ con
        \[
            \langle f,g\rangle = \int_0^1 f(x)g(x)\,dx,
        \]
        para $f(x)=x$ y $g(x)=1-x^2$.
    \end{enumerate}
\end{ejer}

\begin{proof}[Solución]\hspace{0pt}
    \begin{enumerate}
        \item Usando el producto usual
        \[
            \langle x,y\rangle = (1)(4)+(-2)(1)+(0)(-5)+(3)(2)=8.
        \]
        \item Usando $\langle A,B\rangle=\tr(AB^\intercal)$:
        \[
            B^\intercal=\begin{pmatrix}3&4\\-1&2\end{pmatrix},
            \qquad
            AB^\intercal=
            \begin{pmatrix}1&2\\-1&0\end{pmatrix}
            \begin{pmatrix}3&4\\-1&2\end{pmatrix}
            =\begin{pmatrix}1&8\\-3&-4\end{pmatrix}.
        \]
        Por tanto,
        \[
            \langle A,B\rangle=\tr(AB^\intercal)=1+(-4)=-3.
        \]
        \item Tenemos
        \[
            \langle p,q\rangle = 2\cdot 1 + (-1)\cdot 4 + 3\cdot(-1) = -5.
        \]
        \item Calculando la integral:
        \[
            \langle f,g\rangle=\int_0^1 x(1-x^2)\,dx
            =\int_0^1(x-x^3)\,dx
            =\left.\left(\frac{x^2}{2}-\frac{x^4}{4}\right)\right|_0^1
            =\frac14.
        \]
    \end{enumerate}
\end{proof}

%%%%%%%%%%%%%%%%%%%%%%%%%%%%%%%%%%%%%%%%
\begin{ejer}
    Determine si los pares de vectores son ortogonales en el espacio con producto interno dado.
    \begin{enumerate}
        \item En $\R^3$ (producto interno usual):
        \[
            u=(1,-1,2),\quad v=(2,1,0).
        \]
        \item En $\R_3[x]$ con
        \[
            \langle p,q\rangle = a_0b_0+a_1b_1+a_2b_2+a_3b_3,
        \]
        para
        \[
            p(x)=1+2x-x^2,\quad q(x)=2-x+2x^2.
        \]
        \item En $\mathcal{C}([-1,1])$ con
        \[
            \langle f,g\rangle = \int_{-1}^1 f(x)g(x)\,dx,
        \]
        para $f(x)=x$ y $g(x)=1+x^2$.
        \item En $\mathcal{C}([0,\pi])$ con
        \[
            \langle f,g\rangle = \int_0^{\pi} f(x)g(x)\,dx,
        \]
        para $f(x)=\sen x$ y $g(x)=\cos x$.
    \end{enumerate}
\end{ejer}

\begin{proof}[Solución]\hspace{0pt}
    \begin{enumerate}
        \item Usando el producto usual
        \[
            \langle u,v\rangle=(1)(2)+(-1)(1)+(2)(0)=1\neq 0.
        \]
        No son ortogonales.
        \item Calculando el producto interno:
        \[
            \langle p,q\rangle = 1\cdot2 + 2\cdot(-1) + (-1)\cdot2 + 0\cdot0 = -2\neq 0.
        \]
        No son ortogonales.
        \item Calculando el producto interno:
        \[
            \langle f,g\rangle=\int_{-1}^1 x(1+x^2)\,dx=\int_{-1}^1(x+x^3)\,dx=0,
        \]
        porque el integrando es impar en un intervalo simétrico. Son ortogonales.
        \item Calculando el producto interno:
        \[
            \langle f,g\rangle=\int_0^{\pi}\sen x\cos x\,dx
            =\frac12\int_0^{\pi}\sen(2x)\,dx
            =0.
        \]
        Son ortogonales.
    \end{enumerate}
\end{proof}

%%%%%%%%%%%%%%%%%%%%%%%%%%%%%%%%%%%%%%%%
\begin{ejer}
    Sea $W=\gen\{(1,1,0),(1,0,1)\}\subseteq \R^3$ con producto interno usual.
    \begin{enumerate}
        \item Determine una base de $W^\perp$.
        \item Para $v=(2,-1,3)$, calcule $\proy_W(v)$ y $\proy_{W^\perp}(v)$.
    \end{enumerate}
\end{ejer}

\begin{proof}[Solución]\hspace{0pt}
    Sea $a=(1,1,0)$ y $b=(1,0,1)$.
    \begin{enumerate}
        \item Si $x=(x_1,x_2,x_3)\in W^\perp$, entonces
        \[
            x\cdot a=x_1+x_2=0,
            \qquad
            x\cdot b=x_1+x_3=0.
        \]
        De aquí, $x_2=-x_1$ y $x_3=-x_1$, luego
        \[
            x=t(1,-1,-1),\quad t\in\R.
        \]
        Por tanto,
        \[
            W^\perp=\gen\{(1,-1,-1)\}.
        \]
        \item Sea $n=(1,-1,-1)$, base de $W^\perp$. Para $v=(2,-1,3)$:
        \[
            \proy_{W^\perp}(v)=\frac{v\cdot n}{n\cdot n}\,n
            =\frac{2+1-3}{3}n=(0,0,0).
        \]
        Luego,
        \[
            \proy_W(v)=v-\proy_{W^\perp}(v)=(2,-1,3).
        \]
    \end{enumerate}
\end{proof}

%%%%%%%%%%%%%%%%%%%%%%%%%%%%%%%%%%%%%%%%
\begin{ejer}
    Sea $W=\gen\{(1,2,2)\}\subseteq \R^3$ con producto interno usual y sea $v=(3,0,1)$.
    \begin{enumerate}
        \item Encuentre el vector de $W$ más cercano a $v$.
        \item Calcule la distancia de $v$ a $W$.
        \item Verifique el resultado usando el teorema de la proyección.
    \end{enumerate}
\end{ejer}

\begin{proof}[Solución]\hspace{0pt}
    Sea $u=(1,2,2)$, de modo que $W=\gen\{u\}$.
    \begin{enumerate}
        \item El vector de $W$ más cercano a $v$ es $\proy_W(v)$:
        \[
            \proy_W(v)=\frac{v\cdot u}{u\cdot u}\,u
            =\frac{5}{9}(1,2,2)
            =\left(\frac59,\frac{10}9,\frac{10}9\right).
        \]
        \item La distancia es
        \[
            d(v,W)=\|v-\proy_W(v)\|
            =\left\|\left(\frac{22}9,-\frac{10}9,-\frac19\right)\right\|.
        \]
        Entonces
        \[
            d(v,W)=\sqrt{\frac{484+100+1}{81}}=\sqrt{\frac{65}{9}}=\frac{\sqrt{65}}{3}.
        \]
        \item Si $r=v-\proy_W(v)$, se cumple
        \[
            r\cdot u=\frac{22}9-2\cdot\frac{10}9-2\cdot\frac19=0,
        \]
        así que $r\in W^\perp$. Por el teorema de la proyección,
        \[
            v=\proy_W(v)+r,
        \]
        y por tanto $\proy_W(v)$ es el punto de $W$ más cercano a $v$.
    \end{enumerate}
\end{proof}

%%%%%%%%%%%%%%%%%%%%%%%%%%%%%%%%%%%%%%%%
\begin{ejer}
    En un espacio con producto interno, suponga que $u,v\in E$ son ortogonales.
    \begin{enumerate}
        \item Demuestre que $\|u+v\|^2=\|u\|^2+\|v\|^2$.
        \item Aplique el resultado para $u=(2,-1,0)$ y $v=(1,2,0)$ en $\R^3$.
    \end{enumerate}
\end{ejer}

\begin{proof}[Solución]\hspace{0pt}
    \begin{enumerate}
        \item Si $\langle u,v\rangle=0$, entonces
        \[
            \|u+v\|^2=\langle u+v,u+v\rangle
            =\langle u,u\rangle+\langle u,v\rangle+\langle v,u\rangle+\langle v,v\rangle
            =\|u\|^2+\|v\|^2.
        \]
        \item Para $u=(2,-1,0)$ y $v=(1,2,0)$:
        \[
            u\cdot v=2-2+0=0,
        \]
        luego son ortogonales. Además,
        \[
            \|u+v\|^2=\|(3,1,0)\|^2=10,
            \qquad
            \|u\|^2+\|v\|^2=5+5=10.
        \]
        Se verifica Pitágoras.
    \end{enumerate}
\end{proof}


\end{document}